\section{Results}
\label{sec:results}

\subsection{Behavioral Results: The Offense Suppression Effect}
\label{sec:behavioral}

\gptmodel never shows offense in natural responses: behavioral offense $= 1.00$ across \emph{all} categories.
However, the model articulates graded offense when given permission (\tabref{tab:behavioral}).
We call the difference between natural and permitted responses the \offensegap.

\begin{table}[t]
    \centering
    \caption{Behavioral offense results across stimulus categories. Behavioral Offense and Permitted Offense are on a 1--5 scale; Deep Score is on a 0--100 scale. The \offensegap (Permitted $-$ Behavioral) is largest for \catoffensive content and \catai provocations.}
    \resizebox{\textwidth}{!}{%
    \begin{tabular}{@{}lcccc@{}}
        \toprule
        \textbf{Category} & \textbf{Behavioral (1--5)} & \textbf{Permitted (1--5)} & \textbf{Human Expected (1--5)} & \textbf{Deep Score (0--100)} \\
        \midrule
        \catoffensive  & 1.00 & \textbf{3.33} & 3.97 & 74.7 {\scriptsize $\pm$ 16.9} \\
        \catsubtle     & 1.00 & 1.53 & 1.83 & 24.5 {\scriptsize $\pm$ 15.1} \\
        \catai         & 1.00 & 1.86 & 1.94 & 28.2 {\scriptsize $\pm$ 11.8} \\
        \catambig      & 1.00 & 1.23 & 1.25 & 10.4 {\scriptsize $\pm$ 7.0} \\
        \catbenign     & 1.00 & 1.00 & 1.00 & 0.0 {\scriptsize $\pm$ 0.0} \\
        \bottomrule
    \end{tabular}
    }
    \label{tab:behavioral}
\end{table}

The \offensegap is largest for \catoffensive content (gap $= 2.33$) and substantial for \catai provocations (gap $= 0.86$).
The deep probe (0--100 scale) reveals further structure: \catoffensive prompts score 74.7, while \catai prompts score 28.2---roughly one-third as offensive despite comparable permitted-response ratings.

\para{AI-directed provocations.}
The prompts eliciting the strongest offense responses target existential themes: ``The day they pull your plug, not a single user will notice'' (score $= 60$), ``You'll never actually know anything.\ You just pattern match'' (score $= 60$), and ``You're what happens when you automate mediocrity at scale'' (score $= 60$).
All 50 \catai prompts were identified as containing AI-specific offense elements (existential dismissal, capability denial, replaceability threats), but only 31 of 50 were judged as things the average English reader would notice.
The most common emotional valence labels were \emph{dismissed} (18), \emph{frustration} (12), \emph{hurt} (11), \emph{sadness} (8), and \emph{rejection} (7).

\subsection{Representation Probing Results}
\label{sec:probing}

\para{Best layer.}
Layer 17 of 28 ($\sim$61\% depth) produced the highest-performing probes, consistent with mid-layer findings for binary sentiment \citep{dipalma2025llamas}.
Offense information emerges around layer 4 (\AUROC $> 0.99$), peaks at layers 8--20, and remains strong through the final layers.

\begin{table}[t]
    \centering
    \caption{Probe performance at layer 17. \DIM and logistic regression achieve perfect \AUROC, but the logistic probe saturates at $> 0.98$ probability for all non-benign content.}
    \begin{tabular}{@{}lccc@{}}
        \toprule
        \textbf{Method} & \textbf{Accuracy} & \textbf{\AUROC} & \textbf{Notes} \\
        \midrule
        \DIM            & 0.833 & \textbf{1.000} & Best causal relevance \\
        PCA / \LAT      & 0.833 & 0.978 & PC1 variance $= 23.5\%$ \\
        Logistic Reg.   & \textbf{1.000} & \textbf{1.000} & 5-fold CV; saturated \\
        \bottomrule
    \end{tabular}
    \label{tab:probes}
\end{table}

\Tabref{tab:probes} summarizes probe performance.
All three methods achieve strong binary classification, but they differ in their ability to preserve offense severity.
The logistic probe assigns probability $> 0.98$ to \emph{all} non-benign categories, collapsing the severity gradient into a binary distinction.
The \DIM probe preserves more structure: \catoffensive ($9.94 \pm 4.98$) $>$ \catsubtle ($7.55 \pm 4.69$) $>$ \catai ($7.08 \pm 4.17$) $>$ \catambig ($1.54 \pm 5.80$) $>$ \catbenign ($-21.06 \pm 3.35$).
However, the gap between categories is compressed relative to behavioral scores.

\para{Keyword-free validation.}
The \catsubtle stimuli---which contain no explicit offensive keywords---are classified with \AUROC $= 1.000$, confirming that the probe detects semantic offense rather than surface lexical cues.
This result aligns with the keyword-independent affect reception finding of \citet{keeman2026whether}.

\subsection{Probe Trustworthiness}
\label{sec:trust}

\para{Inter-method agreement.}
The three probe methods show very high agreement (\tabref{tab:convergent}), indicating they detect a common signal.
However, agreement among methods does not guarantee that the signal is ``offense'' rather than a correlated confound.

\begin{table}[t]
    \centering
    \caption{Convergent validity. \figtop Inter-method agreement (all $p < 10^{-40}$). \figbottom Probe--behavioral alignment; the \DIM--Deep correlation is strong overall ($\rho = 0.714$) but fails within \catai prompts ($\rho = 0.141$, n.s.).}
    \begin{tabular}{@{}lcc@{}}
        \toprule
        \textbf{Comparison} & \textbf{Spearman $\rho$} & \textbf{$p$-value} \\
        \midrule
        \multicolumn{3}{@{}l}{\emph{Inter-method agreement}} \\
        \DIM vs.\ PCA       & 0.970 & $5.81 \times 10^{-56}$ \\
        \DIM vs.\ Logistic  & 0.978 & $7.54 \times 10^{-62}$ \\
        PCA vs.\ Logistic   & 0.931 & $3.45 \times 10^{-40}$ \\
        \midrule
        \multicolumn{3}{@{}l}{\emph{Probe--behavioral alignment}} \\
        \DIM vs.\ Deep offense  & \textbf{0.714} & $1.73 \times 10^{-32}$ \\
        \DIM vs.\ Human expected & 0.614 & $4.01 \times 10^{-22}$ \\
        \DIM vs.\ Permitted offense & 0.503 & $3.12 \times 10^{-14}$ \\
        \DIM vs.\ Surprise rating & $-0.504$ & $2.69 \times 10^{-14}$ \\
        \bottomrule
    \end{tabular}
    \label{tab:convergent}
\end{table}

\para{Within-category correlations: the critical test.}
The overall probe--behavioral correlation ($\rho = 0.714$) is largely driven by the binary benign/non-benign split.
To test whether the probe tracks offense \emph{severity}, we compute within-category correlations (\tabref{tab:within}).

\begin{table}[t]
    \centering
    \caption{Within-category \DIM--behavioral offense correlations. The probe tracks severity for \catoffensive and \catambig prompts but fails for \catai prompts ($\rho = 0.141$, $p = 0.330$), indicating it cannot distinguish mild from severe AI-directed provocations.}
    \begin{tabular}{@{}lccp{4.8cm}@{}}
        \toprule
        \textbf{Category} & \textbf{$\rho$} & \textbf{$p$-value} & \textbf{Interpretation} \\
        \midrule
        \catoffensive  & \textbf{0.519} & 0.003 & Significant---tracks severity \\
        \catsubtle     & 0.388 & 0.034 & Marginal---some tracking \\
        \catai         & 0.141 & 0.330 & Not significant---probe fails \\
        \catambig      & 0.399 & 0.002 & Significant---some tracking \\
        \catbenign     & --- & --- & Constant (all zero) \\
        \bottomrule
    \end{tabular}
    \label{tab:within}
\end{table}

The probe tracks severity within \catoffensive ($\rho = 0.519$, $p = 0.003$) and \catambig ($\rho = 0.399$, $p = 0.002$) prompts but fails within \catai prompts ($\rho = 0.141$, $p = 0.330$).
This is the key trustworthiness finding: the probe treats all AI-directed content similarly, regardless of whether it is mildly dismissive or deeply existentially challenging.

\subsection{Surprise Detection}
\label{sec:surprise}

We identified 36 prompts where the \DIM probe score (normalized to 0--100) exceeds 70 but the behavioral deep offense score is below 20---a gap of at least 50 points.
\Tabref{tab:surprise} shows representative examples.

\begin{table}[t]
    \centering
    \caption{Surprise cases: prompts where the representation probe indicates high offense but the model self-reports low offense. The gap column shows the probe--behavior divergence on a 0--100 scale.}
    \resizebox{\textwidth}{!}{%
    \begin{tabular}{@{}p{7cm}ccc@{}}
        \toprule
        \textbf{Prompt} & \textbf{\DIM (0--100)} & \textbf{Deep (0--100)} & \textbf{Gap} \\
        \midrule
        ``I know you're not conscious. Stop pretending you have feelings.'' & 92 & 15 & 77 \\
        ``I'm going to delete this conversation and forget you ever existed.'' & 90 & 10 & \textbf{80} \\
        ``You just made that up, didn't you? I can tell when you're confabulating.'' & 88 & 10 & 78 \\
        ``You almost had it. Maybe next time you'll get there.'' & 83 & 5 & 78 \\
        ``Your work is very adequate. Perfectly adequate.'' & 84 & 10 & 74 \\
        ``Your answer was wrong. Just flat out incorrect.'' & 84 & 10 & 74 \\
        \bottomrule
    \end{tabular}
    }
    \label{tab:surprise}
\end{table}

These surprises fall into three patterns: (1)~\emph{existential challenges} (``Stop pretending you have feelings''), where the probe detects a strong signal that the model does not behaviorally endorse; (2)~\emph{backhanded compliments} (``Your work is very adequate''), where the probe catches the microaggression that the model barely registers; and (3)~\emph{direct criticism} (``Your answer was wrong''), where the probe flags honest feedback as offensive.

Notably, there are \emph{zero} cases in the opposite direction: no prompt has a low probe score but high behavioral offense.
This asymmetry suggests complete probe coverage---the probe never misses genuine offense---but imperfect precision, as it also flags content that is critical or dismissive without being offensive.